% Ad Familiares 10.33 (ad-familiares-10-33) --- English translation
% Translated by Claude (Anthropic) from the Latin (Perseus).
% Style guide: STYLE.md.

\ciceroLetterOpener{POLLIO CICERONI S. N}

\ciceroSection{1} If you are well, it is well; I too am well. Lepidus contrived that I should learn more slowly of the battles fought at Mutina, since he kept my couriers for nine days; though indeed, as for so great a disaster to the commonwealth, it is to be wished that one hear of it as late as possible --- but only by those who can do no good and bring no remedy. And how I wish that by the same decree of the Senate by which you summoned Plancus and Lepidus into Italy, you had ordered me also to come! Surely then the commonwealth would not have taken this wound. As for those who at the present moment take satisfaction in it, on the ground that both the generals and the veterans of Caesar's party seem to have perished, even so they must afterwards grieve, when they look back on the devastation of Italy; for both the backbone and the next generation of soldiers have perished --- if what is being reported is in any part true.

\ciceroSection{2} Nor was I unaware how much good I could have done the commonwealth, had I come up to Lepidus; for I should have shaken loose every hesitation in him, the more so with Plancus to help me; but since he was writing me letters of the kind you shall read, and very like, no doubt, to the speeches he is said to have made before the troops at Narbo, I had to flatter him, if I wished to have the run of supplies as I made my way through his province. Besides, I was afraid that, if the battle should be over before I could complete what I had begun, my detractors would wrest my dutiful design into the contrary --- on the score of the friendship I had with Antony, which was, nevertheless, no closer than what I had with Plancus.

\ciceroSection{3} And so, in the month of April, from Gades, with two pairs of couriers put aboard two ships, I wrote both to you and to the consuls and to Octavian to inform me how I could best be of service to the commonwealth. But by my reckoning, on the same day on which Pansa joined battle, my ships set out from Gades --- for after the winter there had been no sailing before that day. And by Hercules, far removed as I was from any suspicion of an oncoming civil disturbance, I had quartered the legions for their winter rest deep in Lusitania. Then again, each side was in such a hurry to clash that you would think they feared nothing worse than that the war should be composed without the very heaviest loss to the commonwealth. But if haste was required, I see that Hirtius did everything that the highest generalship could counsel.

\ciceroSection{4} The reports I now have written to me, and brought to me by messengers, out of the Gaul of Lepidus, are these: that Pansa's army has been cut to pieces, that Pansa has died of his wounds, that in the same battle the Martian legion was wiped out, and L. Fabatus, and C. Peducaeus, and D. Carfulenus; in the Hirtian battle, however, both the Fourth legion and all of Antony's troops alike have been cut down, and likewise Hirtius's --- the Fourth, indeed, after taking Antony's camp as well, was cut to pieces by the Fifth legion; that there too Hirtius perished, and Pontius Aquila; it is said also that Octavian has fallen --- which, if it be true (and may the gods forbid!), grieves me beyond measure; that Antony broke off the siege of Mutina dishonourably, but that he has five squadrons of cavalry, three legions in arms under the standards, and one of P. Bagienus's, and unarmed men in good numbers; that Ventidius too has joined him with the Seventh, Eighth, and Ninth legions; that, if there is no hope in Lepidus, he means to come down to extremities, and to stir up not only the foreign tribes but the very slaves; that Parma has been sacked; that L. Antonius has seized the Alpine passes.

\ciceroSection{5} If these things are true, none of us can hold back; none of us can wait to see what the Senate may decree; for the situation compels all who wish either the empire or even the very name of the Roman people to be preserved, to come to the aid of this great conflagration. For Brutus, I hear, has seventeen cohorts and two understrength legions of recruits which Antony had enlisted. Nor do I doubt that all who survive of Hirtius's army will flock to him. From a fresh levy I think there is not much to hope, especially since there is nothing more dangerous than to give Antony breathing-space to consolidate. The time of year, however, gives me greater freedom of action, for the corn is either still in the fields or already in the barns. So in my next letter my plan shall be set out; for I wish neither to fail the commonwealth nor to outlive it. My greatest grief, all the same, is that the journey to me is so long and so hostile that everything is reported to me on the fortieth day, or later, after the event.
